% Part:sets-functions-relations
% Chapter: size-of-sets
% Section: non-enumerability-alt
% Hindi translation (hi-Deva-IN) of Open Logic commit 9620cc73f9c8e0ad003c514a5d3748f29611c4c0.

\documentclass[../../../include/open-logic-section]{subfiles}

\begin{document}

\olfileid[hi]{sfr}{siz}{nen-alt}

\olsection{\printtoken{S}{nonenumerable} समुच्चय}

\begin{editorial}
  यहाँ \olref[enm-alt]{sec} की परिभाषाओं का उपयोग करके $\Bin^\omega$ और
  $\Pow{\Nat}$ की अगणनीयता सिद्ध करता है; अर्थात~$\Nat$ के साथ एकैकी
  आच्छादक प्रतिचित्रण अपेक्षित है, $\PosInt$ से किसी आच्छादक प्रतिचित्रण के बजाय।
\end{editorial}

\begin{explain}
प्राकृत संख्याओं का समुच्चय $\Nat$ अनंत है। वह स्पष्टतः !!{enumerable} भी
है। लेकिन उल्लेखनीय तथ्य यह है कि \emph{!!{nonenumerable}} समुच्चय भी होते
हैं; अर्थात ऐसे समुच्चय जो !!{enumerable} नहीं हैं
(\olref[sfr][siz][enm-alt]{defn:enumerable} देखें)।

यह आश्चर्यजनक लग सकता है। आखिर यह कहना कि $A$ !!{nonenumerable} है, यह कहना
है कि कोई !!{bijection} $f
\colon \Nat \to A$ \emph{नहीं} है; अर्थात
~$\Nat$ में !!{element} अनंततः अनेक हैं। उन्हें~$A$ में भेजने वाला कोई फलन संपूर्ण~$A$ पर
आच्छादक नहीं होता। अतः यदि $A$ !!{nonenumerable} है, तो~$A$ में प्राकृत
संख्याओं से ``अधिक'' !!{element} हैं।

किसी समुच्चय को !!{nonenumerable} सिद्ध करने के लिए दिखाना होता है कि कोई
उपयुक्त !!{bijection} हो ही नहीं सकता। इसका सर्वोत्तम ढंग यह दिखाना है कि
~$A$ के !!{element} गिनने का प्रत्येक प्रयास कम-से-कम एक !!{element}
अवश्य छोड़ देता है; इससे दिखता है कि कोई फलन $f\colon \Nat
\to A$ !!{surjective} नहीं है। इसे स्थापित करने की सामान्य युक्ति कैंटर की
\emph{विकर्ण विधि} का उपयोग है। मान लीजिए~$A$ के !!{element} $x_1$, $x_2$,
\dots{} के रूप में सूचीबद्ध हैं; तब हम~$A$ का एक और !!{element} बनाते हैं जो
अपने निर्माण के कारण उस सूची में हो ही नहीं सकता।

पर यह सब उदाहरण से सबसे अच्छी तरह समझ आता है। इसलिए हमारा पहला उदाहरण
समुच्चय~$\Bin^\omega$ है, जो $0$ और $1$ की सभी अनंत लड़ियों का समुच्चय है। (`$\Bin$'
द्विआधारी का संकेत है; हम इसे केवल द्वि-अवयवी समुच्चय
$\{0,1\}$ मान सकते हैं।)\oliflabeldef{sfr:card-arithmetic:card-opps:sec}{\footnote{अधिक
ठीक रूप में हमें कहना चाहिए कि $\Bin^\omega$ सभी
$\omega$-अनुक्रमों का समुच्चय है, जो $0$ और $1$ से बने हैं; अर्थात समुच्चय
$\funfromto{\omega}{\{0,1\}}$। लेकिन इसका अर्थ
\olref[sfr][card-arithmetic][card-opps]{sec} में ही स्पष्ट होगा।}{} अभी यह
थोड़ा ढीला निरूपण किसी भ्रम का कारण नहीं होना चाहिए।}
\end{explain}

\begin{thm}
\ollabel{thm:nonenum-bin-omega}
$\Bin^\omega$~!!{nonenumerable} है।
\end{thm}

\begin{proof}
$\Bin^\omega$ के किसी उपसमुच्चय के किसी भी गणन पर विचार कीजिए। अतः हमारे
पास कोई सूची $s_{0}$, $s_{1}$, $s_{2}$, \dots{} है, जहाँ प्रत्येक $s_n$,
$0$ और~$1$ की अनंत लड़ी है। $s_n(m)$ इस सूची की $n$वीं लड़ी का $m$वाँ अंक
हो। तब हम अपनी सूची को ऐसी सरणी मान सकते हैं जिसमें $s_n(m)$ को $n$वीं
पंक्ति और $m$वें स्तंभ में रखा गया है:
\[
\begin{array}{c|c|c|c|c|c}
& 0 & 1 & 2 & 3 & \dots \\\hline
0 & \mathbf{s_{0}(0)} & s_{0}(1) & s_{0}(2) & s_0(3) & \dots \\\hline
1 & s_{1}(0)& \mathbf{s_{1}(1)} & s_1(2) & s_1(3) & \dots \\\hline
2 & s_{2}(0)& s_{2}(1) & \mathbf{s_2(2)} & s_2(3) & \dots \\\hline
3 & s_{3}(0)& s_{3}(1) & s_3(2) & \mathbf{s_3(3)} & \dots \\\hline
\vdots & \vdots & \vdots & \vdots & \vdots & \mathbf{\ddots}
\end{array}
\]
अब हम ऐसी अनंत लड़ी~$d$, जो $0$ और $1$ से बनी है, बनाएँगे जो इस सूची में नहीं है। यह
हम उसके प्रत्येक पद को निर्दिष्ट करके करेंगे; अर्थात हम $d(n)$ बताएँगे और ऐसा
प्रत्येक $n \in \Nat$ के लिए करेंगे। सहज रूप में हम ऊपर की सरणी के विकर्ण
पर नीचे की ओर पढ़ते हैं (इसीलिए इसका नाम ``विकर्ण विधि'' है), फिर प्रत्येक
$1$ को $0$ और प्रत्येक $0$ को~$1$ कर देते हैं। अधिक अमूर्त रूप में हम
$d(n)$ को $0$ या $1$ इस अनुसार परिभाषित करते हैं कि विकर्ण का
$n$वाँ !!{element}, $s_n(n)$, क्रमशः $1$ है या $0$; अर्थात:
\[
d(n) =
\begin{cases}
1 & \text{यदि $s_{n}(n) = 0$}\\
0 & \text{यदि $s_{n}(n) = 1$}
\end{cases}
\]
स्पष्ट है कि $d \in \Bin^\omega$, क्योंकि वह $0$ और $1$ की अनंत लड़ी है।
लेकिन हमने $d$ को इस प्रकार बनाया है कि $d(n) \neq s_n(n)$ किसी भी
$n \in \Nat$ के लिए। अर्थात $d$,~$s_n$ से उसके $n$वें पद पर भिन्न है। अतः
$d \neq s_n$ किसी भी $n\in \Nat$ के लिए। इसलिए $d$ सूची $s_0$, $s_1$,
$s_2$, \dots{} में नहीं हो सकता।

हमने दिखाया कि $\Bin^\omega$ के किसी उपसमुच्चय का कोई भी गणन दिया हो, वह
$\Bin^\omega$ का कोई !!{element} अवश्य छोड़ देता है। इसलिए समुच्चय
$\Bin^\omega$ का कोई गणन नहीं है; अर्थात $\Bin^\omega$
!!{nonenumerable} है।
\end{proof}

\begin{explain}
यह प्रमाण-विधि ``विकर्णीकरण'' कहलाती है, क्योंकि वह~$d$ को परिभाषित करने
के लिए सरणी के विकर्ण का उपयोग करती है। लेकिन विकर्णीकरण में सरणी का होना
आवश्यक नहीं है। वास्तव में, जब कोई सरणी या वास्तविक विकर्ण न हो, तब भी इसी
विचार से किसी समुच्चय को !!{nonenumerable} दिखाया जा सकता है। अगला परिणाम
इसका ढंग स्पष्ट करता है।
\end{explain}

\begin{thm}
\ollabel{thm:nonenum-pownat}
$\Pow{\Nat}$ !!{enumerable} नहीं है।
\end{thm}

\begin{proof}
हम इसी प्रकार आगे बढ़ते हैं: दिखाते हैं कि~$\Nat$ के उपसमुच्चयों की प्रत्येक
सूची $\Nat$ का कोई उपसमुच्चय छोड़ देती है। मान लीजिए हमारे पास कोई सूची
$N_0, N_1, N_2, \ldots$ है, जो $\Nat$ के उपसमुच्चयों की है। हम समुच्चय $D$ इस
प्रकार परिभाषित करते हैं: $n \in D$ यदि और केवल यदि $n \notin N_{n}$:
\[
D = \Setabs{n \in \Nat}{n \notin N_n}
\]
स्पष्ट है कि $D\subseteq \Nat$। लेकिन $D$ सूची में नहीं हो सकता। निर्माण
से $n \in D$ यदि और केवल यदि $n\notin N_n$; अतः $D \neq N_n$ किसी भी
$n \in \Nat$ के लिए।
\end{proof}

\begin{explain}
पिछले प्रमाण में विकर्ण का उल्लेख नहीं था। फिर भी, उसे इस प्रकार चित्रित
करने पर उसमें विकर्ण देखा जा सकता है: कल्पना कीजिए कि समुच्चय $N_0$, $N_1$,
\dots{} किसी सरणी में लिखे हैं, जहाँ~$N_n$ को $n$वीं पंक्ति में इस प्रकार
लिखते हैं कि $m$वें स्तंभ में $m$ तभी लिखते हैं जब $m \in N_n$। उदाहरण के
लिए, मान लीजिए उस सूची के पहले चार समुच्चय $\{0,1,2,\dots\}$,
$\{1, 3, 5, \dots\}$, $\{0,1,4\}$ और $\{2,3,4,\dots\}$ हैं; तब हमारी सरणी
का आरंभ इस प्रकार होगा:
\[
\begin{array}{r@{}rrrrrrr}
  N_0 = \{ & \mathbf{0}, & 1, & 2, & & & & \dots\}\\
  N_1 = \{ &  & \mathbf{1}, &  & 3, &  & 5, & \dots\}\\
  N_2 = \{ & 0, & 1, &  &  & 4\phantom{,} &  & \}\\
  N_3 = \{ &  &  & 2, & \mathbf{3}, & 4, & & \dots\}\\
  &\vdots & & & & & & \ddots\phantom{\}}
  \end{array}
\]
तब $D$ वह समुच्चय है जो विकर्ण पर नीचे की ओर जाकर यह शर्त लगाने से मिलता
है कि $n
\in D$ ठीक तभी, जब $n$ विकर्ण पर \emph{नहीं} हो। इसलिए ऊपर के मामले में हम $0$ और
$1$ को छोड़ेंगे,~$2$ को सम्मिलित करेंगे,~$3$ को छोड़ेंगे, और इसी प्रकार आगे।
\end{explain}

\begin{prob}
स्पष्ट विकर्ण तर्क द्वारा दिखाइए कि सभी फलनों $f \colon \Nat \to \Nat$ का
समुच्चय !!{nonenumerable} है। अर्थात दिखाइए कि यदि $f_1$, $f_2$, \dots{}
फलनों की सूची है और प्रत्येक $f_i\colon
\Nat \to \Nat$ है, तो कोई $g \colon \Nat \to
\Nat$ ऐसा है जो इस सूची में नहीं है।
\end{prob}

\end{document}
